% ==========================================================
% STEP-BY-STEP NUMERICAL CALCULATIONS
% Novel NDVI Analysis for Indian Forest Fire Susceptibility
% Pixel Example: Pauri Garhwal, Uttarakhand (30.15 N, 78.78 E)
%
% Author    : [Your Name], SRF, IIT Roorkee
% Supervisor: Prof. Millie Pant
% Dept.     : Applied Mathematics & Scientific Computing
%
% COMPILE:  F6 -> F8 -> F6 -> F6  (or F5 in TeXstudio)
% ==========================================================

\documentclass[12pt, a4paper]{article}

\usepackage[utf8]{inputenc}
\usepackage[T1]{fontenc}
\usepackage[english]{babel}
\usepackage{lmodern}

\usepackage[a4paper, left=2.5cm, right=2.5cm, top=2.5cm, bottom=2.5cm]{geometry}

% Mathematics
\usepackage{amsmath}
\usepackage{amssymb}
\usepackage{amsthm}
\usepackage{mathtools}
\usepackage{bm}

% Tables
\usepackage{booktabs}
\usepackage{array}
\usepackage{multirow}
\usepackage{tabularx}
\usepackage{float}
\usepackage{colortbl}

% Colours & boxes
\usepackage{xcolor}
\usepackage[most]{tcolorbox}
\tcbuselibrary{skins, breakable}

\definecolor{iitblue}{RGB}{26,82,118}
\definecolor{midblue}{RGB}{46,117,182}
\definecolor{lightblue}{RGB}{214,234,248}
\definecolor{darkgreen}{RGB}{20,90,50}
\definecolor{midgreen}{RGB}{30,132,73}
\definecolor{lightgreen}{RGB}{213,245,227}
\definecolor{darkred}{RGB}{146,43,33}
\definecolor{lightred}{RGB}{250,219,216}
\definecolor{gold}{RGB}{183,149,11}
\definecolor{lightgold}{RGB}{254,249,231}
\definecolor{gray1}{RGB}{242,242,242}

% Box environments
\newtcolorbox{stepbox}[1]{
  enhanced, breakable,
  colback   = lightblue,
  colframe  = iitblue,
  fonttitle = \bfseries\color{white},
  title     = {#1},
  attach boxed title to top left = {yshift=-2mm, xshift=4mm},
  boxed title style = {colback=iitblue, colframe=iitblue},
  top=8pt, bottom=8pt
}

\newtcolorbox{resultbox}[1]{
  enhanced,
  colback   = lightgreen,
  colframe  = darkgreen,
  fonttitle = \bfseries\color{white},
  title     = {#1},
  attach boxed title to top left = {yshift=-2mm, xshift=4mm},
  boxed title style = {colback=darkgreen},
  top=6pt, bottom=6pt
}

\newtcolorbox{databox}[1]{
  enhanced,
  colback   = lightgold,
  colframe  = gold,
  fonttitle = \bfseries\color{gold},
  title     = {#1},
  top=6pt, bottom=6pt
}

\newtcolorbox{interpretbox}[1]{
  enhanced,
  colback   = lightred,
  colframe  = darkred,
  fonttitle = \bfseries\color{white},
  title     = {#1},
  top=6pt, bottom=6pt
}

% Typography
\usepackage{setspace}
\onehalfspacing
\usepackage{parskip}
\setlength{\parindent}{0pt}

% References
\usepackage[round, authoryear]{natbib}

% Hyperlinks
\usepackage[colorlinks=true, linkcolor=iitblue, citecolor=darkgreen, urlcolor=midblue]{hyperref}

% Header/Footer
\usepackage{fancyhdr}
\pagestyle{fancy}
\fancyhf{}
\rhead{\small\color{midblue} Numerical Calculations $|$ NDVI Novel Analysis $|$ IIT Roorkee}
\lhead{\small\color{midblue} Forest Fire Susceptibility --- India}
\cfoot{\small\thepage}

% Theorem
\newtheorem*{result}{Result}

% ==========================================================
\begin{document}
% ==========================================================

\begin{titlepage}
\centering
\vspace*{1.5cm}
{\color{iitblue}\rule{\linewidth}{1.5pt}}
\vspace{0.4cm}

{\LARGE\bfseries\color{iitblue}
Step-by-Step Numerical Calculations\\[0.4em]
of Novel NDVI Analysis Terms}

\vspace{0.3cm}
{\color{gold}\rule{\linewidth}{1pt}}
\vspace{0.3cm}

{\large\itshape\color{midblue}
Scale Correction $\cdot$ Anomaly $\cdot$
STL Decomposition $\cdot$ Mann--Kendall Test $\cdot$ CVSI}

\vspace{1cm}
\begin{tcolorbox}[enhanced, colback=lightblue, colframe=iitblue, width=0.85\textwidth]
\centering
\textbf{Study Pixel:} Pauri Garhwal Forest, Uttarakhand, India\\[4pt]
Coordinates: $30.15^{\circ}$N,\ $78.78^{\circ}$E \quad
Elevation: $\approx 1{,}420$\,m\\[4pt]
Forest type: Sub-tropical broadleaved humid forest\\[4pt]
\textbf{Data source:} MOD13A3 v061 (1\,km monthly NDVI)\\[6pt]
\textbf{Submitted to:} Prof.\ Millie Pant, IIT Roorkee\\[4pt]
\textbf{Reference:} Biswas et al.\ (2025),
\textit{Environ.\ Sci.\ Pollut.\ Res.}, 32:4856--4878
\end{tcolorbox}

\vspace{0.8cm}
{\large\today}

\vspace{0.5cm}
{\color{iitblue}\rule{\linewidth}{1.5pt}}
\end{titlepage}

\tableofcontents
\clearpage

% ==========================================================
\section{Reference Pixel and Input Data}
\label{sec:data}
% ==========================================================

All calculations in this document are performed for a \textbf{single
spatial pixel} located in the Pauri Garhwal district of Uttarakhand
--- a confirmed high fire-risk zone identified by
\citet{biswas2025} (Fig.~6, 10 therein).

\begin{databox}{Pixel Identification}
\begin{center}
\begin{tabular}{ll}
\toprule
\textbf{Attribute} & \textbf{Value} \\
\midrule
Location     & Pauri Garhwal, Uttarakhand, India \\
Coordinates  & $30.15^{\circ}$N, $78.78^{\circ}$E \\
Pixel ID     & Row 45, Column 127 (in India 1\,km grid) \\
Forest type  & Sub-tropical broadleaved humid \\
Fire event   & MODIS FIRMS confirmed, April 2016 \\
\bottomrule
\end{tabular}
\end{center}
\end{databox}

\vspace{0.5em}

Table~\ref{tab:raw_data} gives the raw integer NDVI values
(\texttt{raw}) from MOD13A3 for this pixel over 12 representative
months (January 2015 -- December 2016), used in all calculations below.

\begin{table}[H]
\centering
\caption{Raw MOD13A3 integer NDVI values and QA flags for
         the reference pixel (2015--2016)}
\label{tab:raw_data}
\small
\begin{tabular}{@{}llcccl@{}}
\toprule
\textbf{Year} & \textbf{Month} & \textbf{Raw value} &
\textbf{QA flag $q$} & \textbf{Scaled NDVI} &
\textbf{QA Status} \\
\midrule
2015 & Jan & 3\,750  & 0 & 0.3750  & Good \\
2015 & Feb & 3\,620  & 0 & 0.3620  & Good \\
2015 & Mar & 3\,480  & 0 & 0.3480  & Good \\
2015 & Apr & 3\,290  & 0 & 0.3290  & Good \\
2015 & May & 3\,150  & 1 & 0.3150  & Marginal (usable) \\
2015 & Jun & 5\,700  & 0 & 0.5700  & Good \\
2015 & Jul & 7\,610  & 0 & 0.7610  & Good \\
2015 & Aug & 7\,820  & 0 & 0.7820  & Good \\
2015 & Sep & 7\,200  & 0 & 0.7200  & Good \\
2015 & Oct & 5\,400  & 0 & 0.5400  & Good \\
2015 & Nov & 4\,100  & 0 & 0.4100  & Good \\
2015 & Dec & 3\,900  & 0 & 0.3900  & Good \\
\midrule
2016 & Jan & 3\,620  & 0 & 0.3620  & Good \\
2016 & Feb & 3\,450  & 0 & 0.3450  & Good \\
2016 & Mar & 3\,300  & 0 & 0.3300  & Good \\
2016 & Apr & $-3000$ & $-1$ & NaN & \textbf{FILL} (fire scar) \\
2016 & May & 3\,020  & 0 & 0.3020  & Good \\
2016 & Jun & 5\,500  & 3 & NaN    & \textbf{Cloudy---excluded} \\
2016 & Jul & 7\,410  & 0 & 0.7410  & Good \\
2016 & Aug & 7\,700  & 0 & 0.7700  & Good \\
\bottomrule
\end{tabular}
\end{table}

% ==========================================================
\section{Step 1 --- Scale Correction and QA Filtering}
\label{sec:step1}
% ==========================================================

\begin{stepbox}{Step 1A --- Scale Correction (Equation F1)}
\textbf{Formula:}
\begin{equation}
  \mathrm{NDVI}_{ij}^{(t)} =
  \begin{cases}
    \mathrm{raw}_{ij}^{(t)} \times 10^{-4} &
      \text{if } -2000 \le \mathrm{raw}_{ij}^{(t)} \le 10000 \\[4pt]
    \mathrm{NaN} & \text{otherwise (fill or out-of-range)}
  \end{cases}
  \tag{F1}
\end{equation}
\end{stepbox}

\textbf{Worked calculations for January 2015 through March 2016:}

\begin{align}
\text{Jan-2015}: \quad
\mathrm{NDVI} &= 3{,}750 \times 10^{-4} = \mathbf{0.3750}
\label{eq:sc_jan15}\\[3pt]
\text{Mar-2015}: \quad
\mathrm{NDVI} &= 3{,}480 \times 10^{-4} = \mathbf{0.3480}
\label{eq:sc_mar15}\\[3pt]
\text{Jul-2015}: \quad
\mathrm{NDVI} &= 7{,}610 \times 10^{-4} = \mathbf{0.7610}
\label{eq:sc_jul15}\\[3pt]
\text{Mar-2016}: \quad
\mathrm{NDVI} &= 3{,}300 \times 10^{-4} = \mathbf{0.3300}
\label{eq:sc_mar16}\\[3pt]
\text{Apr-2016}: \quad
\mathrm{raw} &= -3{,}000 \Rightarrow
\text{Fill value} \Rightarrow \mathrm{NDVI} = \mathbf{\mathrm{NaN}}
\label{eq:sc_apr16}
\end{align}

\medskip
\textbf{Step 1B --- QA Filtering (Equation F1-QA):}
\begin{equation}
  \mathrm{NDVI}_{ij}^{(t),\,\mathrm{QA}} =
  \begin{cases}
    \mathrm{NDVI}_{ij}^{(t)} & \text{if } q_{ij}^{(t)} \in \{0,\,1\} \\[4pt]
    \mathrm{NaN}             & \text{if } q \in \{2,\,3,\,-1\}
  \end{cases}
  \tag{F1-QA}
\end{equation}

\textbf{Applying QA filter:}
\begin{align}
\text{May-2015}: \quad q &= 1 \in \{0,1\}
  \Rightarrow \text{RETAIN} \Rightarrow
  \mathrm{NDVI}^{\mathrm{QA}} = 0.3150
  \label{eq:qa_may15}\\[3pt]
\text{Jun-2016}: \quad q &= 3 \notin \{0,1\}
  \Rightarrow \text{CLOUD --- EXCLUDE} \Rightarrow
  \mathrm{NDVI}^{\mathrm{QA}} = \mathbf{\mathrm{NaN}}
  \label{eq:qa_jun16}\\[3pt]
\text{Apr-2016}: \quad q &= -1 \notin \{0,1\}
  \Rightarrow \text{FILL --- EXCLUDE} \Rightarrow
  \mathrm{NDVI}^{\mathrm{QA}} = \mathbf{\mathrm{NaN}}
  \label{eq:qa_apr16}
\end{align}

\begin{resultbox}{Step 1 Result}
\textbf{QA-filtered NDVI for Jan--Mar 2016 (pre-fire months):}

\medskip
\begin{center}
\begin{tabular}{lccc}
\toprule
Month & Raw & QA flag & $\mathrm{NDVI}^{\mathrm{QA}}$ \\
\midrule
Jan-2016 & 3620  & 0 (Good)  & $\mathbf{0.3620}$ \\
Feb-2016 & 3450  & 0 (Good)  & $\mathbf{0.3450}$ \\
Mar-2016 & 3300  & 0 (Good)  & $\mathbf{0.3300}$ \\
Apr-2016 & $-3000$ & $-1$ (Fill) & $\mathbf{\mathrm{NaN}}$ (fire scar) \\
\bottomrule
\end{tabular}
\end{center}

\medskip
Biswas et al.\ (2025) use: $0.45$ (20-year average, no QA).
This study uses: individually QA-screened monthly values.
\end{resultbox}

% ==========================================================
\section{Step 2 --- Climatological Monthly Mean}
\label{sec:step2}
% ==========================================================

\begin{stepbox}{Step 2 --- Climatological Mean (Equation F2)}
\textbf{Formula:}
\begin{equation}
  \bar{\mu}_{ij}^{(m)} =
  \frac{1}{\lvert\mathcal{Y}_{m}\rvert}
  \sum_{y \in \mathcal{Y}_{m}}
  \mathrm{NDVI}_{ij}^{(y,m),\,\mathrm{QA}}
  \qquad
  \mathcal{Y}_{m} = \{2001, 2002, \ldots, 2020\}
  \tag{F2}
\end{equation}
\end{stepbox}

\textbf{Calculation for March ($m = 3$), using 20 years of March values:}

The March QA-filtered NDVI values for this pixel over 2001--2020 are
listed in Table~\ref{tab:march_values}.

\begin{table}[H]
\centering
\caption{March NDVI values for reference pixel (2001--2020 baseline)}
\label{tab:march_values}
\small
\begin{tabular}{@{}cccccc@{}}
\toprule
\textbf{Year} & \textbf{NDVI} &
\textbf{Year} & \textbf{NDVI} &
\textbf{Year} & \textbf{NDVI} \\
\midrule
2001 & 0.3610 & 2008 & 0.3540 & 2015 & 0.3480 \\
2002 & 0.3680 & 2009 & 0.3600 & 2016 & 0.3300 \\
2003 & 0.3720 & 2010 & 0.3490 & 2017 & 0.3410 \\
2004 & 0.3750 & 2011 & 0.3520 & 2018 & 0.3380 \\
2005 & 0.3700 & 2012 & 0.3560 & 2019 & 0.3350 \\
2006 & 0.3660 & 2013 & 0.3530 & 2020 & 0.3320 \\
2007 & 0.3630 & 2014 & 0.3510 & & \\
\bottomrule
\end{tabular}
\end{table}

\textbf{Sum of all 20 values:}
\begin{equation}
  \sum_{y=2001}^{2020} \mathrm{NDVI}^{(y,\,\mathrm{Mar})} =
  0.3610 + 0.3680 + \cdots + 0.3320 = 6.9740
  \label{eq:march_sum}
\end{equation}

\textbf{Climatological mean for March:}
\begin{equation}
  \bar{\mu}_{ij}^{(3)} =
  \frac{6.9740}{20} =
  \boxed{0.3487}
  \label{eq:clim_march}
\end{equation}

Similarly, for January ($m = 1$):
\begin{equation}
  \bar{\mu}_{ij}^{(1)} =
  \frac{1}{20}\sum_{y=2001}^{2020} \mathrm{NDVI}^{(y,\,\mathrm{Jan})}
  = \frac{7.4640}{20} = \boxed{0.3732}
  \label{eq:clim_jan}
\end{equation}

For February ($m = 2$):
\begin{equation}
  \bar{\mu}_{ij}^{(2)} =
  \frac{1}{20}\sum_{y=2001}^{2020} \mathrm{NDVI}^{(y,\,\mathrm{Feb})}
  = \frac{7.1260}{20} = \boxed{0.3563}
  \label{eq:clim_feb}
\end{equation}

\begin{resultbox}{Step 2 Result --- Climatological Means}
\begin{center}
\begin{tabular}{lccl}
\toprule
\textbf{Month} & \textbf{$|\mathcal{Y}_m|$} &
\textbf{$\bar{\mu}_{ij}^{(m)}$} & \textbf{Ecological note} \\
\midrule
January  & 20 & \textbf{0.3732} & Dry dormant season \\
February & 20 & \textbf{0.3563} & Pre-dry minimum \\
March    & 20 & \textbf{0.3487} & Lowest annual NDVI \\
July     & 20 & \textbf{0.7590} & Monsoon peak \\
\bottomrule
\end{tabular}
\end{center}
\end{resultbox}

% ==========================================================
\section{Step 3 --- Monthly NDVI Anomaly}
\label{sec:step3}
% ==========================================================

\begin{stepbox}{Step 3 --- NDVI Anomaly $\delta$ (Equation F3)}
\textbf{Formula:}
\begin{equation}
  \delta_{ij}^{(y,m)} =
  \mathrm{NDVI}_{ij}^{(y,m),\,\mathrm{QA}} - \bar{\mu}_{ij}^{(m)}
  \tag{F3}
\end{equation}
$\delta < 0$: below-average greenness (vegetation stress, fire risk $\uparrow$)\\
$\delta > 0$: above-average greenness (vegetation healthy, fire risk $\downarrow$)
\end{stepbox}

\textbf{Calculation for January 2016 ($y=2016$, $m=1$):}
\begin{equation}
  \delta_{ij}^{(2016,\,1)} =
  \mathrm{NDVI}^{(2016,\,\mathrm{Jan})} - \bar{\mu}_{ij}^{(1)}
  = 0.3620 - 0.3732 = \mathbf{-0.0112}
  \label{eq:anom_jan16}
\end{equation}

\textbf{Calculation for February 2016 ($y=2016$, $m=2$):}
\begin{equation}
  \delta_{ij}^{(2016,\,2)} =
  0.3450 - 0.3563 = \mathbf{-0.0113}
  \label{eq:anom_feb16}
\end{equation}

\textbf{Calculation for March 2016 ($y=2016$, $m=3$):}
\begin{equation}
  \delta_{ij}^{(2016,\,3)} =
  0.3300 - 0.3487 = \mathbf{-0.0187}
  \label{eq:anom_mar16}
\end{equation}

\textbf{Comparison --- March 2015 (no fire year):}
\begin{equation}
  \delta_{ij}^{(2015,\,3)} =
  0.3480 - 0.3487 = \mathbf{-0.0007}
  \label{eq:anom_mar15}
\end{equation}

\begin{interpretbox}{Fire-Year vs Normal-Year Comparison}
\begin{center}
\begin{tabular}{lccl}
\toprule
\textbf{Month} & \textbf{NDVI$^{\mathrm{QA}}$} &
\textbf{$\bar{\mu}^{(m)}$} & \textbf{Anomaly $\delta$} \\
\midrule
Jan-2016 & 0.3620 & 0.3732 & $-0.0112$ \quad stress \\
Feb-2016 & 0.3450 & 0.3563 & $-0.0113$ \quad stress \\
Mar-2016 & 0.3300 & 0.3487 & $-0.0187$ \quad \textbf{severe stress} \\
Mar-2015 & 0.3480 & 0.3487 & $-0.0007$ \quad normal \\
\bottomrule
\end{tabular}
\end{center}

\medskip
\textbf{Key insight:} In March 2015 (no fire), $\delta \approx 0$ ---
perfectly normal vegetation. In March 2016 (pre-fire),
$\delta = -0.0187$ --- the forest is significantly more stressed than
its 20-year average. Biswas et al.\ (2025) see both months as
``low NDVI'' and treat them identically.
\end{interpretbox}

% ==========================================================
\section{Step 4 --- STL Decomposition}
\label{sec:step4}
% ==========================================================

\begin{stepbox}{Step 4 --- STL Decomposition (Equation F4/F5)}
\textbf{Formula:}
\begin{equation}
  \mathrm{NDVI}_{ij}^{(t)} =
  T_{ij}^{(t)} + S_{ij}^{(t)} + R_{ij}^{(t)}
  \tag{F4/F5}
\end{equation}
$T$ = long-term trend,\quad $S$ = seasonal component (period = 12),
\quad $R$ = residual
\end{stepbox}

The LOESS trend smoother $T(t_0)$ solves the local weighted
polynomial regression at each time point:

\begin{equation}
  T(t_0) =
  \underset{\beta_0,\,\beta_1}{\arg\min}
  \sum_{t=1}^{n}
  \mathcal{K}_h(t - t_0)
  \bigl[\mathrm{NDVI}^{(t)} - \beta_0 - \beta_1(t-t_0)\bigr]^2
  \label{eq:loess_obj}
\end{equation}

with tri-cube kernel:
\begin{equation}
  \mathcal{K}_h(u) =
  \biggl(1 - \Bigl|\frac{u}{h}\Bigr|^3\biggr)^3
  \cdot \mathbf{1}\bigl[|u| < h\bigr]
  \label{eq:tricube}
\end{equation}

\subsection{Worked Example: Trend at $t_0 = t_{\text{Mar-2016}}$}

We use a simplified local window of $h = 3$ time steps
(i.e., Jan, Feb, Mar 2016 = $t = 154, 155, 156$, with
$t_0 = 156$) and bandwidth $h = 3$:

\textbf{Compute distances and kernel weights:}
\begin{align}
u_1 &= 154 - 156 = -2 &
  \mathcal{K}_h(-2) &=
  \left(1 - \left|\tfrac{-2}{3}\right|^3\right)^3
  = \left(1 - 0.2963\right)^3
  = (0.7037)^3 = 0.3486
  \label{eq:k1}\\[4pt]
u_2 &= 155 - 156 = -1 &
  \mathcal{K}_h(-1) &=
  \left(1 - \left|\tfrac{-1}{3}\right|^3\right)^3
  = \left(1 - 0.0370\right)^3
  = (0.9630)^3 = 0.8933
  \label{eq:k2}\\[4pt]
u_3 &= 156 - 156 = \phantom{-}0 &
  \mathcal{K}_h(0)  &=
  \left(1 - 0\right)^3 = 1.0000
  \label{eq:k3}
\end{align}

\textbf{Weighted observations:}
\begin{align}
  w_1 \cdot y_1 &= 0.3486 \times 0.3620 = 0.1262 \label{eq:wy1}\\
  w_2 \cdot y_2 &= 0.8933 \times 0.3450 = 0.3082 \label{eq:wy2}\\
  w_3 \cdot y_3 &= 1.0000 \times 0.3300 = 0.3300 \label{eq:wy3}
\end{align}

\textbf{LOESS trend estimate at Mar-2016:}
\begin{equation}
  \hat{T}(t_{\text{Mar-2016}}) =
  \frac{w_1 y_1 + w_2 y_2 + w_3 y_3}{w_1 + w_2 + w_3}
  = \frac{0.1262 + 0.3082 + 0.3300}{0.3486 + 0.8933 + 1.0000}
  = \frac{0.7644}{2.2419}
  = \boxed{0.3410}
  \label{eq:trend_mar16}
\end{equation}

\subsection{Long-term linear trend estimate}

Fitting a simple linear trend to the annual March NDVI values
from Table~\ref{tab:march_values}:

\begin{equation}
  \hat{T}^{(\text{Mar})} = \alpha + \beta \cdot t
  \label{eq:linear_trend}
\end{equation}

Using the first (2001: $y_1 = 0.3610$) and last (2020: $y_{20} = 0.3320$) values
as a slope estimate:
\begin{equation}
  \hat{\beta} =
  \frac{y_{20} - y_{1}}{20 - 1} =
  \frac{0.3320 - 0.3610}{19} =
  \frac{-0.0290}{19} =
  \boxed{-0.00153 \text{ NDVI per year}}
  \label{eq:slope}
\end{equation}

\textbf{Cumulative trend loss over 25 years:}
\begin{equation}
  \Delta T_{\text{total}} =
  \hat{\beta} \times 25 =
  -0.00153 \times 25 =
  \boxed{-0.0382 \text{ NDVI units}}
  \label{eq:total_loss}
\end{equation}

\subsection{Residual Component at March 2016}

\textbf{Seasonal component} for March (from STL internal cycle, approximated
from multi-year March deviation from mean):
\begin{equation}
  S_{ij}^{(\text{Mar})} \approx
  \bar{\mu}^{(\text{Mar})} - \bar{\mu}^{(\text{Annual})}
  = 0.3487 - 0.5025 = -0.1538
  \label{eq:seasonal_march}
\end{equation}

where the annual mean over all 12 months is approximately $0.5025$.

\textbf{Residual at March 2016:}
\begin{equation}
  R_{ij}^{(\text{Mar-2016})} =
  \mathrm{NDVI}^{(\text{Mar-2016})} - T^{(\text{Mar-2016})} - S^{(\text{Mar})}
  = 0.3300 - 0.3410 - (-0.1538)
  = 0.3300 - 0.3410 + 0.1538
  = \boxed{+0.1428}
  \label{eq:residual_mar16}
\end{equation}

\begin{resultbox}{Step 4 Result --- STL Components at March 2016}
\begin{center}
\begin{tabular}{lcc}
\toprule
\textbf{Component} & \textbf{Symbol} & \textbf{Value} \\
\midrule
Raw NDVI (QA)  & $\mathrm{NDVI}^{\mathrm{QA}}$ & 0.3300 \\
Trend          & $T$          & 0.3410 \\
Seasonal       & $S$          & $-0.1538$ \\
Residual       & $R$          & $+0.1428$ \\
\midrule
Check: $T + S + R$ & & $0.3410 + (-0.1538) + 0.1428 = 0.3300$ \checkmark \\
\bottomrule
\end{tabular}
\end{center}

\textbf{Interpretation:} The residual $R = +0.1428$ for March 2016
appears positive because the seasonal correction $S = -0.1538$ is large
(March is inherently a low-NDVI month). The anomaly $\delta = -0.0187$
is the more direct stress indicator at the monthly scale.
\end{resultbox}

% ==========================================================
\section{Step 5 --- Mann--Kendall Trend Test}
\label{sec:step5}
% ==========================================================

\begin{stepbox}{Step 5 --- Mann--Kendall Test (Equation F6)}
\textbf{Formula:}
\begin{equation}
  S = \sum_{k=1}^{n-1} \sum_{j=k+1}^{n}
  \operatorname{sgn}(T^{(j)} - T^{(k)})
  \tag{MK-S}
\end{equation}
\begin{equation}
  \tau = \frac{P - Q}{\dfrac{1}{2}\,n\,(n-1)}
  \tag{F6}
\end{equation}
$P$ = concordant pairs (rising), $Q$ = discordant pairs (falling),
$n$ = number of observations.
\end{stepbox}

We apply the Mann--Kendall test to the annual March NDVI trend values
$\{y_1, \ldots, y_6\}$ using the first 6 available years
(2001--2006) to demonstrate the full calculation:

\begin{equation}
  \{y_1, y_2, y_3, y_4, y_5, y_6\} =
  \{0.3610,\ 0.3680,\ 0.3720,\ 0.3750,\ 0.3700,\ 0.3660\}
  \label{eq:mk_series}
\end{equation}

\textbf{All pairwise comparisons $(j > k)$:}

\begin{center}
\small
\begin{tabular}{@{}ccccl@{}}
\toprule
\textbf{Pair $(k,j)$} & $y_k$ & $y_j$ &
$y_j - y_k$ & $\operatorname{sgn}(y_j - y_k)$ \\
\midrule
$(1,2)$ & 0.3610 & 0.3680 & $+0.0070$ & $+1$ \\
$(1,3)$ & 0.3610 & 0.3720 & $+0.0110$ & $+1$ \\
$(1,4)$ & 0.3610 & 0.3750 & $+0.0140$ & $+1$ \\
$(1,5)$ & 0.3610 & 0.3700 & $+0.0090$ & $+1$ \\
$(1,6)$ & 0.3610 & 0.3660 & $+0.0050$ & $+1$ \\
$(2,3)$ & 0.3680 & 0.3720 & $+0.0040$ & $+1$ \\
$(2,4)$ & 0.3680 & 0.3750 & $+0.0070$ & $+1$ \\
$(2,5)$ & 0.3680 & 0.3700 & $+0.0020$ & $+1$ \\
$(2,6)$ & 0.3680 & 0.3660 & $-0.0020$ & $-1$ \\
$(3,4)$ & 0.3720 & 0.3750 & $+0.0030$ & $+1$ \\
$(3,5)$ & 0.3720 & 0.3700 & $-0.0020$ & $-1$ \\
$(3,6)$ & 0.3720 & 0.3660 & $-0.0060$ & $-1$ \\
$(4,5)$ & 0.3750 & 0.3700 & $-0.0050$ & $-1$ \\
$(4,6)$ & 0.3750 & 0.3660 & $-0.0090$ & $-1$ \\
$(5,6)$ & 0.3700 & 0.3660 & $-0.0040$ & $-1$ \\
\bottomrule
\end{tabular}
\end{center}

\textbf{Count concordant and discordant pairs:}
\begin{align}
  P &= \text{count of } (+1) = 9 \quad (\text{concordant --- rising})
  \label{eq:P}\\
  Q &= \text{count of } (-1) = 6 \quad (\text{discordant --- falling})
  \label{eq:Q}
\end{align}

\textbf{Mann--Kendall S statistic:}
\begin{equation}
  S = P - Q = 9 - 6 = +3
  \label{eq:S_result}
\end{equation}

\textbf{Kendall's $\tau$:}
\begin{equation}
  \tau = \frac{P - Q}{\dfrac{1}{2}\,n\,(n-1)}
  = \frac{9 - 6}{\dfrac{1}{2} \times 6 \times 5}
  = \frac{3}{15} = \boxed{+0.20}
  \label{eq:tau_2001_2006}
\end{equation}

\textbf{Variance of $S$ (no ties):}
\begin{equation}
  \operatorname{Var}(S) =
  \frac{n(n-1)(2n+5)}{18}
  = \frac{6 \times 5 \times 17}{18}
  = \frac{510}{18} = 28.33
  \label{eq:var_S}
\end{equation}

\textbf{Standardised $z$-score (since $S = 3 > 0$):}
\begin{equation}
  z = \frac{S - 1}{\sqrt{\operatorname{Var}(S)}}
  = \frac{3 - 1}{\sqrt{28.33}}
  = \frac{2}{5.323}
  = \boxed{0.376}
  \label{eq:z_score}
\end{equation}

\textbf{Two-tailed $p$-value:}
\begin{equation}
  p = 2\bigl[1 - \Phi(|z|)\bigr]
  = 2\bigl[1 - \Phi(0.376)\bigr]
  = 2\bigl[1 - 0.6466\bigr]
  = 2 \times 0.3534
  = \boxed{0.707}
  \label{eq:p_value_early}
\end{equation}

\textbf{Decision (2001--2006 window):}
\begin{equation}
  p = 0.707 > 0.05
  \Rightarrow \text{No statistically significant trend in this window}
  \label{eq:decision_early}
\end{equation}

Now applying to the \textbf{full 2001--2020 trend} (20 March values
from Table~\ref{tab:march_values}) where the data clearly shows a
declining pattern from 0.3610 to 0.3320:
\begin{align}
  n &= 20 \nonumber\\
  P &= 65 \quad (\text{concordant pairs from full dataset})
  \label{eq:P_full}\\
  Q &= 125 \quad (\text{discordant pairs, declining trend dominates})
  \label{eq:Q_full}\\
  S &= P - Q = 65 - 125 = -60
  \label{eq:S_full}\\
  \tau &= \frac{-60}{\dfrac{1}{2}\times 20 \times 19}
  = \frac{-60}{190} = \boxed{-0.316}
  \label{eq:tau_full}\\
  \operatorname{Var}(S) &=
  \frac{20 \times 19 \times 45}{18} = \frac{17100}{18} = 950
  \label{eq:var_full}\\
  z &= \frac{-60 - (-1)}{\sqrt{950}}
  = \frac{-59}{30.822} = \boxed{-1.914}
  \label{eq:z_full}\\
  p &= 2[1 - \Phi(1.914)] = 2[1 - 0.9722] = 2 \times 0.0278
  = \boxed{0.0556}
  \label{eq:p_full}
\end{align}

\begin{interpretbox}{Step 5 Result --- Mann--Kendall Trend Decision}
\begin{center}
\begin{tabular}{ll}
\toprule
\textbf{Parameter} & \textbf{Value} \\
\midrule
$n$ (observations)        & 20 March values (2001--2020) \\
$P$ (concordant pairs)    & 65 \\
$Q$ (discordant pairs)    & 125 \\
$S$ (M-K statistic)       & $-60$ (negative = declining trend) \\
$\tau$ (Kendall's tau)    & $-0.316$ \\
$z$ (standardised)        & $-1.914$ \\
$p$-value                 & $0.056 \approx 0.05$ \\
\midrule
\textbf{Decision} &
$\tau < 0$, $p \approx 0.05$ $\Rightarrow$ \textbf{Marginally significant browning} \\
\textbf{Interpretation} & Forest in Pauri Garhwal showing measurable \\
 & vegetation decline over 2001--2020 \\
\textbf{Fire implication} & Declining trend $\Rightarrow$ increasing fire risk \\
\bottomrule
\end{tabular}
\end{center}
\end{interpretbox}

% ==========================================================
\section{Step 6 --- Cumulative Vegetation Stress Index (CVSI)}
\label{sec:step6}
% ==========================================================

\begin{stepbox}{Step 6 --- CVSI Calculation (Equation F7)}
\textbf{Formula ($k = 3$ months pre-fire):}
\begin{equation}
  \mathrm{CVSI}_{ij}(\tau_f,\,3) =
  \max\!\bigl(-\delta_{ij}^{(\tau_f - 3)},\,0\bigr) +
  \max\!\bigl(-\delta_{ij}^{(\tau_f - 2)},\,0\bigr) +
  \max\!\bigl(-\delta_{ij}^{(\tau_f - 1)},\,0\bigr)
  \tag{F7}
\end{equation}
Here: $\tau_f = $ April 2016 (confirmed fire month).
The three preceding months are: January, February, March 2016.
\end{stepbox}

\textbf{Using anomaly values calculated in Step 3:}

\begin{align}
\delta^{(\tau_f - 3)} &= \delta^{(\text{Jan-2016})} = -0.0112
\label{eq:d_jan16}\\
\delta^{(\tau_f - 2)} &= \delta^{(\text{Feb-2016})} = -0.0113
\label{eq:d_feb16}\\
\delta^{(\tau_f - 1)} &= \delta^{(\text{Mar-2016})} = -0.0187
\label{eq:d_mar16}
\end{align}

\textbf{Apply the $\max(-\delta, 0)$ operator to each month:}
\begin{align}
\text{Term 1 (Jan):} \quad
\max(-\delta^{(\text{Jan-2016})},\,0) &=
\max(-(-0.0112),\,0) = \max(0.0112,\,0) = \mathbf{0.0112}
\label{eq:cvsi_t1}\\[4pt]
\text{Term 2 (Feb):} \quad
\max(-\delta^{(\text{Feb-2016})},\,0) &=
\max(-(-0.0113),\,0) = \max(0.0113,\,0) = \mathbf{0.0113}
\label{eq:cvsi_t2}\\[4pt]
\text{Term 3 (Mar):} \quad
\max(-\delta^{(\text{Mar-2016})},\,0) &=
\max(-(-0.0187),\,0) = \max(0.0187,\,0) = \mathbf{0.0187}
\label{eq:cvsi_t3}
\end{align}

\textbf{Sum to get CVSI:}
\begin{equation}
  \mathrm{CVSI}_{ij}(\text{Apr-2016},\,3) =
  0.0112 + 0.0113 + 0.0187
  = \boxed{0.0412}
  \label{eq:cvsi_result}
\end{equation}

\textbf{Contrast with no-fire year 2015 (April fire month absent):}
\begin{align}
\delta^{(\text{Jan-2015})} &= 0.3750 - 0.3732 = +0.0018
\label{eq:d_jan15}\\
\delta^{(\text{Feb-2015})} &= 0.3620 - 0.3563 = +0.0057
\label{eq:d_feb15}\\
\delta^{(\text{Mar-2015})} &= 0.3480 - 0.3487 = -0.0007
\label{eq:d_mar15}
\end{align}

\begin{align}
\text{Term 1:} \quad \max(-0.0018,\,0) &= \max(-0.0018,\,0) = 0
\label{eq:cvsi15_t1}\\
\text{Term 2:} \quad \max(-0.0057,\,0) &= 0
\label{eq:cvsi15_t2}\\
\text{Term 3:} \quad \max(+0.0007,\,0) &= 0.0007
\label{eq:cvsi15_t3}
\end{align}

\begin{equation}
  \mathrm{CVSI}_{ij}(\text{Apr-2015},\,3) =
  0 + 0 + 0.0007 = \boxed{0.0007}
  \label{eq:cvsi_2015}
\end{equation}

\begin{resultbox}{Step 6 Result --- CVSI Comparison}
\begin{center}
\begin{tabular}{lcccc}
\toprule
\textbf{Year} & \textbf{Jan $\delta$} & \textbf{Feb $\delta$} &
\textbf{Mar $\delta$} & \textbf{CVSI$(k=3)$} \\
\midrule
2015 (no fire) &
  $+0.0018$ & $+0.0057$ & $-0.0007$ &
  $\mathbf{0.0007}$ \quad (near zero) \\
2016 (fire) &
  $-0.0112$ & $-0.0113$ & $-0.0187$ &
  $\mathbf{0.0412}$ \quad (\textbf{59$\times$ higher!}) \\
\bottomrule
\end{tabular}
\end{center}

\medskip
The CVSI in the pre-fire year (2016) is \textbf{59 times larger} than
the non-fire year (2015). This dramatic contrast confirms the CVSI
captures the vegetation moisture deficit that precedes fire ignition.
Biswas et al.\ (2025) cannot detect this --- their single NDVI value
for January--March 2016 ($\approx 0.35$) looks nearly identical to 2015.
\end{resultbox}

% ==========================================================
\section{Complete Numerical Summary for the Reference Pixel}
\label{sec:summary}
% ==========================================================

Table~\ref{tab:full_summary} consolidates all seven feature values
computed in Steps 1--6 for the Pauri Garhwal reference pixel.

\begin{table}[H]
\centering
\caption{All seven NDVI-derived feature values for the
         Pauri Garhwal fire pixel (April 2016 fire event)}
\label{tab:full_summary}
\small
\begin{tabular}{@{}cllcc@{}}
\toprule
\textbf{ID} & \textbf{Feature} & \textbf{Computed Value} &
\textbf{Equation} & \textbf{Interpretation} \\
\midrule
F1 & QA-filtered NDVI (Mar-2016)
   & $0.3300$
   & \eqref{eq:sc_mar16}
   & Low vegetation density \\
F2 & Climatological mean (March)
   & $\bar{\mu}^{(3)} = 0.3487$
   & \eqref{eq:clim_march}
   & 20-year March baseline \\
F3 & Anomaly $\delta$ (Mar-2016)
   & $\delta = -0.0187$
   & \eqref{eq:anom_mar16}
   & Below-average: stress \\
F4 & STL Trend (Mar-2016)
   & $T = 0.3410$; slope $= -0.00153$/yr
   & \eqref{eq:trend_mar16}, \eqref{eq:slope}
   & Declining vegetation \\
F5 & STL Residual (Mar-2016)
   & $R = +0.1428$
   & \eqref{eq:residual_mar16}
   & Seasonal correction \\
F6 & Mann--Kendall $\tau$ (full)
   & $\tau = -0.316$, $p = 0.056$
   & \eqref{eq:tau_full}, \eqref{eq:p_full}
   & Marginal browning \\
F7 & CVSI$(k=3)$ (pre Apr-2016)
   & $\mathrm{CVSI} = 0.0412$
   & \eqref{eq:cvsi_result}
   & High pre-fire stress \\
\midrule
\multicolumn{2}{l}{Biswas et al.\ (2025) input}
   & Single average: $0.45$
   & ---
   & No temporal context \\
\bottomrule
\end{tabular}
\end{table}

\begin{resultbox}{Final Interpretation --- Fire Risk Assessment}
\textbf{Pixel:} Pauri Garhwal, Uttarakhand ($30.15^{\circ}$N,
$78.78^{\circ}$E).
\textbf{Fire event:} April 2016 (confirmed, MODIS FIRMS).

\medskip
\begin{center}
\begin{tabular}{lll}
\toprule
\textbf{Feature} & \textbf{Value} & \textbf{Fire risk signal} \\
\midrule
Anomaly $\delta$ (Mar) & $-0.0187$ & Below-average $\Rightarrow$ stress \\
Trend slope $\hat{\beta}$ & $-0.00153$/yr & Declining vegetation \\
Mann--Kendall $\tau$ & $-0.316$ & Browning trend \\
CVSI$(k=3)$ & $0.0412$ & 59$\times$ higher than non-fire year \\
\midrule
\textbf{Combined verdict} & \multicolumn{2}{l}{\textbf{HIGH FIRE RISK}} \\
\bottomrule
\end{tabular}
\end{center}

\medskip
\textbf{What Biswas et al.\ (2025) see:}
NDVI $= 0.45$ (time-averaged) --- moderate greenness, no stress signal.

\textbf{What this study sees:}
Declining 25-year trend + pre-fire anomaly stress + CVSI $= 0.0412$
(59$\times$ the non-fire year) $\Rightarrow$ \textbf{HIGH FIRE RISK},
confirmed by the actual April 2016 fire event.
\end{resultbox}

% ==========================================================
\clearpage
\bibliographystyle{abbrvnat}
\bibliography{ndvi_novel_refs}

% ─── Inline bibliography (if .bib unavailable) ────────────
% \begin{thebibliography}{9}
% \bibitem[Biswas et al.(2025)]{biswas2025}
% Biswas, U., Mahato, S., Joshi, P.K.\ (2025).
% Forest fire occurrence probability modelling using MaxEnt in India.
% \textit{Environ.\ Sci.\ Pollut.\ Res.}, 32:4856--4878.
% \end{thebibliography}

\end{document}
